\section{Introduction}

LLM assistants have become a routine part of terminal work. A typical user, an ML
researcher training models across a laptop and a handful of remote GPU machines,
babysits long runs in a shell, edits configs, holds sessions open with a multiplexer,
and asks an assistant to explain a CUDA out-of-memory trace or draft an \code{rsync}
command. In practice the assistant is the weakest link, for two reasons that compound. It is generic: it runs in a separate window,
sees only the fragment the user pastes into it, and does not know that in this project a
training run is launched with a particular \code{sbatch} invocation settled on last
week. It is also amnesiac: it forgets the command the user corrected yesterday, forgets ``where
was I'' on reattach, and knows nothing about the editor
it sits beside. In the remote setting that defines GPU work there is a third problem:
reaching the work at all usually means exporting an API key onto a machine the user does
not own, where it lingers in the environment, the shell history, and often a dotfile.

The gap is not a shortage of AI tooling. IDEs have capable agents \citep{cursor}, and
per-invocation CLI tools bring code generation to the shell
\citep{anthropic2024claudecode,copilotcli}; but no modern agentic tool sits
\emph{inside} the terminal session and empowers the user there --- close to the
action, where the command, the log, and the stack trace already are, with nothing
pasted to and fro. AI terminals such as Warp \citep{warp} transform the terminal
application itself, so they cannot follow the user onto a headless server, and they
see only the commands run inside them; classic multiplexers \citep{tmux,zellij}
persist but carry no intelligence, and history tools \citep{atuin,thefuck} treat the
user's history as valuable but never feed it to a model. MARS bridges the gap: a
terminal-native application that runs anywhere a shell runs, including inside an ssh
session on the box where the work is.

Our observation is that an agent embedded at the terminal has direct access to memory
that a browser-tab assistant never sees: the user's own history of commands they
actually ran, and the tool's own documentation, action registry, and configuration
schema. Neither corpus is
large, and neither requires training. What they require is an agent that lives beside
them, persists between requests, and captures corrections as they happen. CLI agents
can read the same local files MARS reads; however, at
the time of writing their persistent context is curated by hand (project memory files)
rather than captured automatically from the commands a user accepts, and every remote
machine needs its own credential. MARS differs in the combination: one persistent
process that captures every accepted correction, grounds the agent in its own action
registry, and keeps the credential home. We built
it to satisfy all three conditions\footnote{Code, gold sets, and frozen result files:
\url{https://github.com/anjishnu/mars} (MIT). Install:
\code{cargo install mars-terminal}
(\url{https://crates.io/crates/mars-terminal}). Screencast: [URL].}; it hosts the
user's shell in its panes rather than replacing it, retrieves from both corpora with
plain lexical search, and injects the results into the prompt.

We expected local memory to help. We did not expect a method this simple to nearly close
the gap between a small model and its personalized ceiling, even when the stored memory
is only a \emph{paraphrase} of the user's earlier request. On a 24-item set of
project-specific commands, retrieving a handful of reworded past corrections lifts
\code{claude-haiku-4-5} from 38\% (9/24) to 96\% (23/24) under a protocol in which
a win requires generalizing across a paraphrase gap, not looking up a planted answer. In a narrow domain, memory does not
need to be sophisticated to be decisive.

This paper makes four contributions. Firstly, an open-source, single-binary
memory-augmented terminal agent (\S\ref{sec:agent}), built from minimal parts: a registry-validated directive protocol, a
corrective-memory loop, and a model cascade. Secondly, two kinds of \emph{local} memory
(\S\ref{sec:m1}--\ref{sec:m2}): a \emph{reflex} memory of the user's own commands, and
a \emph{reflective} memory of the tool's own ground truth. Thirdly, the architecture (\S\ref{sec:system}) behind always-present
local memory: a persistent daemon, and an SSH broker that
keeps the API key off remote machines. Finally, a self-instrumenting evaluation
(\S\ref{sec:eval}) of corrective memory under a leakage-controlled protocol. MARS is about 9{,}500 lines of
Rust~\citep{kumar2026mars}.
